%----------------------------------------------------------------------------------------
% VEELTERMFUNCTIES. Alle opgaven zijn eigen werk.
%
% Opbouw: Start, leerstof (theorie, werkwijze of uitleg, voorbeeld, korte oefening), en op een
% nieuwe pagina de Opdrachten (door elkaar, stijgende moeilijkheid), Koppel en Test jezelf.
% Labels van dit hoofdstuk beginnen met h1-.
%----------------------------------------------------------------------------------------

\mychapter{1}{Veeltermfuncties}

\lkr{Leerplan B+S: LPD 17 en K1 (rationale functies), met de veeltermfuncties als voorbereiding: delen, de reststelling, de factorstelling en het schema van Horner, ontbinden en de grafiek schetsen vanuit nulwaarden, multipliciteit en de hoogstegraadsterm. Eerste- en tweedegraadsfuncties kennen leerlingen uit de 2de graad.}

Eerste- en tweedegraadsfuncties ken je al: rechten en parabolen. Een veeltermfunctie van hogere graad kan meer bochten maken. Om haar grafiek te schetsen heb je vooral haar nulwaarden nodig, en daarvoor moet je de veelterm ontbinden. Dat doe je door te delen.

Het hoofdstuk volgt die weg. Eerst leer je veeltermen delen, met een staartdeling en met het snelle schema van Horner. Daarna zie je waarom een nulwaarde en een factor hetzelfde zijn, en hoe je nulwaarden vindt. Op het einde schets je de grafiek met de nulwaarden, hun multipliciteit en de hoogstegraadsterm. In hoofdstuk 2 heb je dat allemaal nodig voor rationale functies.

\begin{oefeningen}[Start]
\oef{h1-start-1}

\oef{h1-verken}
\end{oefeningen}

%========================================================================================
\section{Veeltermen delen}

\subsection{Veeltermen en hun graad}

Een veelterm is een som van termen $a x^k$, met een getal $a$ en een natuurlijke exponent $k$. Je hebt er al veel uitgewerkt en ontbonden. In dit hoofdstuk krijgen de onderdelen een vaste naam, want de stellingen spreken erover.

\begin{theorie}
Een \textbf{\textcolor{accent}{veeltermfunctie}} van \textbf{\textcolor{accent}{graad}} $n$ is

\qquad $\formule{f(x) = a_n x^n + a_{n-1}x^{n-1} + \cdots + a_1 x + a_0}$ \quad met $a_n \neq 0$.

De getallen $a_k$ zijn de \textbf{\textcolor{accent}{coëfficiënten}}, $a_n x^n$ is de \textbf{\textcolor{accent}{hoogstegraadsterm}} en $a_0 = f(0)$ is de \textbf{\textcolor{accent}{constante term}}. De nulveelterm $0$ heeft geen graad.
\end{theorie}

Moet je een product uitwerken om te weten welke graad het heeft? Nee: de hoogste macht ontstaat maar op één manier.

\begin{theorie}
\begin{stelling}{Graad van een product}
Heeft $A(x)$ graad $n$ en hoogstegraadsterm $a_n x^n$, en $B(x)$ graad $m$ en hoogstegraadsterm $b_m x^m$, dan heeft $A(x) \cdot B(x)$

\qquad graad $\formule{n + m}$ \quad en hoogstegraadsterm \quad $\formule{a_n b_m\, x^{n+m}}$
\end{stelling}

\begin{bewijs}
Werk je het product uit, dan is elke term het product van een term van $A$ en een term van $B$:
\bewijsregel{a_i x^i \cdot b_j x^j = a_i b_j\, x^{i+j} \qquad \text{met } i \leqslant n \text{ en } j \leqslant m .}
Dus is $i + j \leqslant n + m$, met gelijkheid alleen als $i = n$ en $j = m$. De enige term in $x^{n+m}$ is dus $a_n b_m x^{n+m}$, en die valt niet weg, want $a_n \neq 0$ en $b_m \neq 0$ geeft $a_n b_m \neq 0$. Andere termen hebben een lagere graad.
\end{bewijs}
\end{theorie}

Je vermenigvuldigt dus alleen de hoogstegraadstermen van de factoren. Voor een macht geldt hetzelfde: $(A(x))^k$ heeft graad $kn$ en hoogstegraadsterm $(a_n x^n)^k$.

\voorbeeld{}{Neem $f(x) = (1 - 2x)^3(x^2 + 4)$. Bepaal de graad, de hoogstegraadsterm en de constante term.\label{vb:h1-graad}}{%
\uitleg{Werk niet uit. De hoogstegraadsterm is het product van de hoogstegraadstermen van de factoren, en de constante term is $f(0)$.}
\begin{stappen}
(-2x)^3 \cdot x^2 = -8x^3 \cdot x^2 = -8x^5 & De hoogste term van $1 - 2x$ is $-2x$, niet $1$. Let op het teken bij de derde macht. \\[1ex]
\text{graad } 5 & De exponent van de hoogstegraadsterm, of $3 \cdot 1 + 2$. \\[1ex]
f(0) = 1^3 \cdot 4 = 4 & De constante term.
\end{stappen}}

Bij een som werkt dat niet zomaar: de hoogste termen kunnen tegen elkaar wegvallen. Zo heeft $(x^2 + 1)^2 - x^4 = 2x^2 + 1$ graad 2 en niet 4.

\begin{oefeningen}
\needspace{6\baselineskip}
\oef{h1-graad}
\end{oefeningen}

\subsection{De euclidische deling}

Deel je $147$ door $12$, dan vind je quotiënt $12$ en rest $3$: \ $147 = 12 \cdot 12 + 3$, met een rest die kleiner is dan de deler. Veeltermen deel je op dezelfde manier. Wat \textit{kleiner} betekent, meet je nu met de graad: de rest heeft een lagere graad dan de deler.

\begin{theorie}
\begin{stelling}{Euclidische deling}
Voor veeltermen $A(x)$ en $B(x) \neq 0$ bestaan er unieke veeltermen $Q(x)$ (het \textbf{\textcolor{accent}{quotiënt}}) en $R(x)$ (de \textbf{\textcolor{accent}{rest}}) met

\qquad $\formule{A(x) = B(x) \cdot Q(x) + R(x)}$ \quad en \quad $\formule{R(x) = 0 \text{ of graad } R < \text{graad } B}$
\end{stelling}

\begin{bewijs}
\textit{Bestaan.} Stel dat $B$ graad $m$ heeft, met hoogstegraadsterm $b_m x^m$. Is $A = 0$ of graad $A < m$, neem dan $Q = 0$ en $R = A$. Heeft $A$ graad $n \geqslant m$ en hoogstegraadsterm $a_n x^n$, trek dan af:
\bewijsregel{A_1(x) = A(x) - \frac{a_n}{b_m}x^{n-m} \cdot B(x) .}
De hoogstegraadsterm van het product is $a_n x^n$, dus die valt weg: $A_1 = 0$ of $A_1$ heeft een lagere graad dan $A$. Herhaal dat met $A_1$, en zo verder. Omdat de graad bij elke stap daalt, kom je na hoogstens $n - m + 1$ stappen bij een veelterm $R$ met $R = 0$ of graad $R < m$. Tel je de termen op die je met $B$ hebt vermenigvuldigd, dan heb je $Q$, en $A = BQ + R$.

\textit{Uniek.} Stel $A = BQ_1 + R_1 = BQ_2 + R_2$, met beide resten $0$ of van graad kleiner dan $m$. Dan is
\bewijsregel{B(x) \cdot \bigl(Q_1(x) - Q_2(x)\bigr) = R_2(x) - R_1(x) .}
Is $Q_1 \neq Q_2$, dan heeft het linkerlid graad minstens $m$ (graad van een product), en het rechterlid is $0$ of heeft graad kleiner dan $m$. Dat kan niet. Dus $Q_1 = Q_2$, en dan ook $R_1 = R_2$.
\end{bewijs}
\end{theorie}

Het bewijs is meteen een rekenmethode: telkens de hoogstegraadsterm wegwerken. Dat is de staartdeling, zoals bij getallen.

\voorbeeld{}{Deel $A(x) = x^4 + 2x^3 - x + 3$ door $B(x) = x^2 + 1$.\label{vb:h1-staartdeling}}{%
\uitleg{Deel telkens de hoogstegraadsterm van wat overblijft door $x^2$, vermenigvuldig terug met $B(x)$ en trek af. Stop als de graad kleiner is dan 2.}
\begin{stappenbreed}
x^4 + 2x^3 - x + 3 - x^2(x^2 + 1) = 2x^3 - x^2 - x + 3 & Eerste term van $Q$: $x^4 : x^2 = x^2$. \\
2x^3 - x^2 - x + 3 - 2x(x^2 + 1) = -x^2 - 3x + 3 & Tweede term: $2x^3 : x^2 = 2x$. \\
-x^2 - 3x + 3 - (-1)(x^2 + 1) = -3x + 4 & Derde term: $-1$. Graad 1: stop. \\
Q(x) = x^2 + 2x - 1,\quad R(x) = -3x + 4 & Tel de termen van $Q$ op.
\end{stappenbreed}
Controle: $(x^2 + 1)(x^2 + 2x - 1) - 3x + 4 = x^4 + 2x^3 - x + 3$.}

\begin{werkwijze}
Orden $A(x)$ en $B(x)$ volgens dalende machten. Deel de hoogstegraadsterm van $A(x)$ door die van $B(x)$: dat is de eerste term van het quotiënt. Trek die term maal $B(x)$ af. Herhaal met wat overblijft, tot de graad kleiner is dan die van $B(x)$: dat is de rest. Controleer met $A = B \cdot Q + R$.
\end{werkwijze}

\begin{oefeningen}
\needspace{6\baselineskip}
\oef{h1-staart}
\end{oefeningen}

\subsection{Het schema van Horner}

Delen door $x - a$ komt zo vaak voor dat er een snellere manier is. Het schema van Horner werkt alleen met de coëfficiënten.

\voorbeeld{}{Deel $A(x) = 2x^3 - 3x^2 + 4x - 5$ door $x - 2$.\label{vb:h1-horner}}{%
\naast{%
\uitleg{Hier is $a = 2$. Zet de coëfficiënten op een rij en breng de eerste naar beneden. Vermenigvuldig die met $2$, zet het resultaat onder de volgende coëfficiënt en tel op. Herhaal.}
\begin{stappenvrij}{4.6cm}{3.6cm}
2 \cdot 2 = 4,\ -3 + 4 = 1 & Tweede kolom. \\
2 \cdot 1 = 2,\ 4 + 2 = 6 & Derde kolom. \\
2 \cdot 6 = 12,\ -5 + 12 = 7 & Laatste kolom: de rest. \\
Q(x) = 2x^2 + x + 6,\ R = 7 & Een graad lager dan $A$.
\end{stappenvrij}}{%
{\renewcommand{\arraystretch}{1.3}%
$\begin{array}{c|rrrr}
 & 2 & -3 & 4 & -5 \\
2 & & 4 & 2 & 12 \\ \hline
 & 2 & 1 & 6 & \textcolor{accent}{\mathbf{7}}
\end{array}$}}}

\begin{werkwijze}
Schrijf de coëfficiënten van $A(x)$ op een rij, volgens dalende machten, met een $0$ voor elke ontbrekende macht. Zet $a$ links: bij $x + 3 = x - (-3)$ is $a = -3$. Breng de eerste coëfficiënt naar beneden, en herhaal: vermenigvuldig met $a$, zet het resultaat in de volgende kolom en tel op. Het laatste getal is de rest, de andere zijn de coëfficiënten van het quotiënt.
\end{werkwijze}

Waarom werkt dat? Neem een veelterm van graad 3. Het quotiënt bij deling door $x - a$ heeft graad 2, en de rest is een getal $r$:
\[a_3x^3 + a_2x^2 + a_1x + a_0 = (x - a)(q_2x^2 + q_1x + q_0) + r = q_2x^3 + (q_1 - aq_2)x^2 + (q_0 - aq_1)x + (r - aq_0) .\]
\naast{Twee veeltermen zijn gelijk als hun coëfficiënten gelijk zijn. Dat geeft $q_2 = a_3$, en daarna telkens: de vorige maal $a$, plus de volgende coëfficiënt:
\[q_1 = a_2 + aq_2, \qquad q_0 = a_1 + aq_1, \qquad r = a_0 + aq_0 .\]
Dat is precies wat het schema rechts doet. Bij een hogere graad gaat het net zo.}{%
{\renewcommand{\arraystretch}{1.3}%
$\begin{array}{c|cccc}
 & a_3 & a_2 & a_1 & a_0 \\
a & & aq_2 & aq_1 & aq_0 \\ \hline
 & q_2 & q_1 & q_0 & r
\end{array}$}}

Horner werkt alleen bij een deler $x - a$. Door $x^2 + 1$ deel je met een staartdeling.

\begin{oefeningen}
\needspace{6\baselineskip}
\oef{h1-delen}
\end{oefeningen}

\subsection{De reststelling}

In voorbeeld~\ref{vb:h1-horner} is de rest $7$. Vul je $2$ in, dan krijg je $A(2) = 16 - 12 + 8 - 5 = 7$: hetzelfde getal. In de Verken zag je iets gelijkaardigs. Dat is geen toeval.

\begin{theorie}
\begin{stelling}{Reststelling}
De rest bij deling van $A(x)$ door $x - a$ is \quad $\formule{A(a)}$
\end{stelling}

\begin{bewijs}
De deler $x - a$ heeft graad 1, dus de rest is $0$ of heeft graad 0: in beide gevallen een getal $r$. Uit de euclidische deling
\bewijsregel{A(x) = (x - a) \cdot Q(x) + r}
volgt voor $x = a$:
\bewijsregel{A(a) = 0 \cdot Q(a) + r = r .}
\end{bewijs}
\end{theorie}

Je kent de rest dus zonder te delen: vul de nulwaarde van de deler in. Dat is handig als je alleen de rest nodig hebt, of om een deling met Horner te controleren. Is de deler $2x - 1$, dan werkt hetzelfde bewijs met de nulwaarde $\frac12$: de rest is $A\bigl(\frac12\bigr)$.

\voorbeeld{}{Bepaal de rest bij deling van $A(x) = x^{20} - 3x^5 + 2$ door $x + 1$.\label{vb:h1-rest}}{%
\uitleg{Met Horner heb je een rij van 21 coëfficiënten nodig, bijna allemaal nul. De reststelling geeft de rest in één regel.}
\begin{stappen}
x + 1 = x - (-1) & Dus $a = -1$. \\[1ex]
A(-1) = (-1)^{20} - 3 \cdot (-1)^5 + 2 & Vul $-1$ in. \\[1ex]
= 1 + 3 + 2 = 6 & Een even macht van $-1$ is $1$, een oneven macht is $-1$.
\end{stappen}
De rest is $6$.}

Bij een deler van graad 2 of meer werkt dit niet: de rest is dan geen getal maar een veelterm, zoals $px + q$, en één invulling geeft je die niet.

\begin{oefeningen}
\needspace{6\baselineskip}
\oef{h1-horner}

\needspace{6\baselineskip}
\oef{h1-rest}
\end{oefeningen}

%========================================================================================
\section{Ontbinden en nulwaarden}

Om de nulwaarden van een veelterm van graad 3 of meer te vinden, ontbind je hem in factoren: een product is nul als een factor nul is. Maar hoe vind je een factor? Het verband tussen factoren en nulwaarden is de sleutel.

\subsection{De factorstelling}

\begin{theorie}
$B(x)$ is een \textbf{\textcolor{accent}{deler}} van $A(x)$ als de rest van de deling $0$ is, dus als $\formule{A(x) = B(x) \cdot Q(x)}$ voor een veelterm $Q(x)$.

\begin{stelling}{Factorstelling}
$x - a$ is een deler van $A(x)$ \quad $\formule{\iff}$ \quad $\formule{A(a) = 0}$
\end{stelling}

\begin{bewijs}
Volgens de reststelling is $A(x) = (x - a) \cdot Q(x) + A(a)$.

$\Leftarrow$ \ Is $A(a) = 0$, dan staat er $A(x) = (x - a) \cdot Q(x)$: $x - a$ is een deler.

$\Rightarrow$ \ Is $x - a$ een deler, dan is $A(x) = (x - a) \cdot Q(x)$ voor een veelterm $Q(x)$, en dus
\bewijsregel{A(a) = (a - a) \cdot Q(a) = 0 .}
\end{bewijs}
\end{theorie}

Een nulwaarde vinden en een factor vinden is dus hetzelfde. Is $a$ een nulwaarde, dan deel je door $x - a$ met Horner, en houd je een quotiënt van een graad lager over.

\voorbeeld{}{Bepaal $k$ zodat $x - 2$ een deler is van $A(x) = x^3 + kx^2 - 4x + 12$. Ontbind $A(x)$ dan volledig.\label{vb:h1-factor}}{%
\uitleg{Volgens de factorstelling moet $A(2) = 0$. Dat is een vergelijking in $k$. Daarna deel je door $x - 2$.}
\begin{stappen}
A(2) = 8 + 4k - 8 + 12 = 4k + 12 & Vul $2$ in. \\[1ex]
4k + 12 = 0 \iff k = -3 & De rest moet nul zijn. \\[1ex]
A(x) = (x - 2)(x^2 - x - 6) & Horner met $a = 2$ op $1, -3, -4, 12$ geeft $1, -1, -6$ en rest $0$. \\[1ex]
A(x) = (x - 2)(x - 3)(x + 2) & Som $1$, product $-6$.
\end{stappen}}

\begin{oefeningen}
\needspace{6\baselineskip}
\oef{h1-factor}
\end{oefeningen}

\subsection{Gehele nulwaarden zoeken}

Om de factorstelling te gebruiken, heb je een nulwaarde nodig. Lukraak getallen proberen is hopeloos. Maar als alle coëfficiënten gehele getallen zijn, is er een korte lijst met kandidaten.

\begin{theorie}
\begin{stelling}{Gehele nulwaarden}
Heeft $A(x) = a_n x^n + \cdots + a_1 x + a_0$ gehele coëfficiënten, dan is elke gehele nulwaarde van $A(x)$ een \textbf{\textcolor{accent}{deler van de constante term}} $\formule{a_0}$.
\end{stelling}

\begin{bewijs}
Stel dat het gehele getal $a$ een nulwaarde is. Dan is
\bewijsregel{a_n a^n + a_{n-1}a^{n-1} + \cdots + a_1 a + a_0 = 0 ,}
waaruit
\bewijsregel{a_0 = -a \cdot \bigl(a_n a^{n-1} + a_{n-1}a^{n-2} + \cdots + a_1\bigr) .}
Tussen de haakjes staat een geheel getal, want alle $a_k$ en $a$ zijn geheel. Dus is $a_0$ een veelvoud van $a$: $a$ is een deler van $a_0$.
\end{bewijs}
\end{theorie}

De stelling zegt niet dat elke deler een nulwaarde is: ze geeft alleen de kandidaten. Die probeer je één voor één, positief en negatief.

\voorbeeld{}{Ontbind $f(x) = x^3 - 2x^2 - 5x + 6$ en geef de nulwaarden.\label{vb:h1-ontbinden}}{%
\naast{%
\uitleg{De kandidaten zijn de delers van $6$: $\pm 1, \pm 2, \pm 3, \pm 6$. Vind je een nulwaarde, deel dan met Horner en ontbind het quotiënt (een tweedegraadsveelterm) verder.}
\begin{stappenvrij}{4.9cm}{3.2cm}
f(1) = 1 - 2 - 5 + 6 = 0 & $x - 1$ is een deler. \\
f(x) = (x - 1)(x^2 - x - 6) & Horner. \\
= (x - 1)(x - 3)(x + 2) & Som $1$, product $-6$. \\
\text{nulwaarden: } -2,\ 1,\ 3 &
\end{stappenvrij}}{%
{\renewcommand{\arraystretch}{1.3}%
$\begin{array}{c|rrrr}
 & 1 & -2 & -5 & 6 \\
1 & & 1 & -1 & -6 \\ \hline
 & 1 & -1 & -6 & \textcolor{accent}{\mathbf{0}}
\end{array}$}}}

\begin{werkwijze}
Zonder eerst een gemeenschappelijke factor af en gebruik merkwaardige producten. Probeer daarna de delers van de constante term (positief en negatief) tot je een nulwaarde $a$ vindt. Deel door $x - a$ met Horner, en ontbind het quotiënt verder. Een tweedegraadsveelterm ontbind je met de discriminant.
\end{werkwijze}

De methode vindt alleen gehele nulwaarden. Zo heeft $2x^3 + x^2 + x - 1$ geen gehele nulwaarde: de kandidaten $1$ en $-1$ geven $3$ en $-3$. Toch is $\frac12$ een nulwaarde, en $2x^3 + x^2 + x - 1 = (2x - 1)(x^2 + x + 1)$. Een grafiek of een rekentoestel helpt je dan op weg; met de factorstelling controleer je exact. Is de constante term $0$, dan zonder je eerst $x$ af: elk getal is een deler van $0$, dus de lijst met kandidaten zegt dan niets.

\begin{oefeningen}
\needspace{6\baselineskip}
\oef{h1-kandidaat}

\needspace{6\baselineskip}
\oef{h1-ontbinden}
\end{oefeningen}

\subsection{Multipliciteit}

Soms komt dezelfde factor meer dan eens voor. Na deling door $x - a$ kan $a$ opnieuw een nulwaarde van het quotiënt zijn.

\begin{theorie}
Is $A(x) = (x - a)^m \cdot Q(x)$ met $Q(a) \neq 0$, dan is $a$ een nulwaarde met \textbf{\textcolor{accent}{multipliciteit}} $m$. Bij $m = 1$ spreek je van een enkelvoudige, bij $m = 2$ van een dubbele nulwaarde.
\end{theorie}

\voorbeeld{}{Ontbind $f(x) = x^3 + x^2 - 8x - 12$ en geef de nulwaarden met hun multipliciteit.\label{vb:h1-multipliciteit}}{%
\uitleg{Probeer de delers van $12$. Blijft de gevonden nulwaarde een nulwaarde van het quotiënt, dan komt de factor twee keer voor.}
\begin{stappenbreed}
f(-2) = -8 + 4 + 16 - 12 = 0 & $x + 2$ is een deler. \\[1ex]
f(x) = (x + 2)(x^2 - x - 6) & Horner met $a = -2$ op $1, 1, -8, -12$ geeft $1, -1, -6$. \\[1ex]
f(x) = (x + 2)(x + 2)(x - 3) = (x + 2)^2(x - 3) & Ook het quotiënt is nul in $-2$. \\[1ex]
-2 \text{ met multipliciteit } 2,\quad 3 \text{ met multipliciteit } 1 & Tel de exponenten: $2 + 1 = 3$, de graad.
\end{stappenbreed}}

\begin{oefeningen}
\needspace{6\baselineskip}
\oef{h1-multipliciteit}
\end{oefeningen}

In het voorbeeld is de som van de multipliciteiten precies de graad. Meer nulwaarden kunnen er nooit zijn.

\begin{theorie}
\begin{stelling}{Aantal nulwaarden}
Een veelterm van graad $n$ heeft \textbf{\textcolor{accent}{hoogstens $n$ nulwaarden}}, geteld met hun multipliciteit.
\end{stelling}

\begin{bewijs}
Neem verschillende nulwaarden $c_1, c_2, \ldots, c_k$ van $A(x)$, met multipliciteit $m_1, m_2, \ldots, m_k$. Dan is $A(x) = (x - c_1)^{m_1}Q_1(x)$. Voor $c_2$ geeft dat
\bewijsregel{0 = A(c_2) = (c_2 - c_1)^{m_1} \cdot Q_1(c_2) ,}
en omdat $c_2 \neq c_1$, is $Q_1(c_2) = 0$. Dus is $c_2$ een nulwaarde van $Q_1$. De factor $(x - c_1)^{m_1}$ is niet nul in $c_2$, dus de hele factor $(x - c_2)^{m_2}$ van $A$ zit in $Q_1$: \ $Q_1(x) = (x - c_2)^{m_2}Q_2(x)$. Zo verder met de andere nulwaarden geeft
\bewijsregel{A(x) = (x - c_1)^{m_1}(x - c_2)^{m_2} \cdots (x - c_k)^{m_k} \cdot Q(x) \qquad \text{met } Q(x) \neq 0 .}
Volgens de graad van een product is $n = m_1 + m_2 + \cdots + m_k + \text{graad } Q$, dus $m_1 + m_2 + \cdots + m_k \leqslant n$.
\end{bewijs}
\end{theorie}

Ken je $n$ nulwaarden (met multipliciteit) van een veelterm van graad $n$, dan ken je hem dus op een constante factor na: $Q(x)$ heeft dan graad $0$. Zo stel je een voorschrift op.

\voorbeeld{}{Bepaal de veeltermfunctie $f$ van graad 3 met een dubbele nulwaarde $-1$, een nulwaarde $2$ en $f(0) = 4$.\label{vb:h1-voorschrift}}{%
\naast{%
\uitleg{De multipliciteiten zijn samen $3$, de graad. Er is dus geen ruimte voor een andere factor dan een constante $c$.}
\begin{stappenvrij}{4.3cm}{3.6cm}
f(x) = c\,(x + 1)^2(x - 2) & Aantal nulwaarden. \\[1ex]
f(0) = c \cdot 1 \cdot (-2) = -2c & Vul $0$ in. \\[1ex]
-2c = 4 \iff c = -2 & \\[1ex]
f(x) = -2(x + 1)^2(x - 2) &
\end{stappenvrij}}{%
\begin{tikzpicture}[x=.62cm,y=.2cm,font=\scriptsize]
  \rooster{-2}{1}{3}{-6}{2}{8}
  \kromme{-1.9}{2.25}{-2*(\x+1)^2*(\x-2)}
  \vol{(-1,0)} \vol{(2,0)} \vol{(0,4)}
\end{tikzpicture}}}

\begin{oefeningen}
\needspace{6\baselineskip}
\oef{h1-voorschrift}
\end{oefeningen}

%========================================================================================
\section{De grafiek van een veeltermfunctie}

Met de nulwaarden weet je waar de grafiek de $x$-as ontmoet. Om haar te schetsen moet je nog weten wat ze daar doet (snijden of raken) en waar ze naartoe gaat aan de uiteinden.

\subsection{Het gedrag aan de uiteinden}

Voor grote $|x|$ overheerst de hoogstegraadsterm: in $f(x) = x^3 - 2x^2 - 5x + 6$ is voor $x = 100$ de term $x^3 = 1\,000\,000$, en de rest samen maar ongeveer $-20\,500$. De grafiek gedraagt zich dan zoals die van $a_n x^n$.

\begin{theorie}
\begin{stelling}{Gedrag aan de uiteinden}
Voor grote $|x|$ heeft een veeltermfunctie van graad $n \geqslant 1$ hetzelfde teken als haar \textbf{\textcolor{accent}{hoogstegraadsterm}} $\formule{a_n x^n}$, en worden de functiewaarden onbeperkt groot in absolute waarde.
\end{stelling}

\begin{bewijs}
Zonder de hoogstegraadsterm af:
\bewijsregel{f(x) = a_n x^n \Bigl(1 + \frac{a_{n-1}}{a_n x} + \frac{a_{n-2}}{a_n x^2} + \cdots + \frac{a_0}{a_n x^n}\Bigr) .}
Voor grote $|x|$ wordt elke breuk tussen de haakjes willekeurig klein, want de teller is een vast getal en de noemer wordt onbeperkt groot. De factor tussen de haakjes ligt dan dicht bij $1$, en is zeker positief. Dus heeft $f(x)$ hetzelfde teken als $a_n x^n$, en is $|f(x)|$ ongeveer $|a_n x^n|$, dat onbeperkt groot wordt.
\end{bewijs}
\end{theorie}

\textit{Willekeurig klein} en \textit{onbeperkt groot} maak je later precies met limieten. Het teken van $a_n x^n$ hangt af van het teken van $a_n$, en van de vraag of $n$ even of oneven is. Dat geeft vier gevallen.

\begin{center}
\renewcommand{\arraystretch}{1.3}
\begin{tabular}{|c|c|c|}
\hline
 & $n$ even & $n$ oneven \\ \hline
$a_n > 0$ & \begin{tabular}{c}\begin{tikzpicture}[x=.55cm,y=.25cm]
  \draw[->,grijs] (-2.4,0) -- (2.5,0); \draw[->,grijs] (0,-4.3) -- (0,4.6);
  \kromme{-2.02}{2.02}{(\x)^4-3*(\x)^2}
\end{tikzpicture}\\ \small $x^4 - 3x^2$: links en rechts boven\end{tabular} & \begin{tabular}{c}\begin{tikzpicture}[x=.55cm,y=.25cm]
  \draw[->,grijs] (-2.4,0) -- (2.5,0); \draw[->,grijs] (0,-4.3) -- (0,4.6);
  \kromme{-2.22}{2.22}{(\x)^3-3*\x}
\end{tikzpicture}\\ \small $x^3 - 3x$: links onder, rechts boven\end{tabular} \\ \hline
$a_n < 0$ & \begin{tabular}{c}\begin{tikzpicture}[x=.55cm,y=.25cm]
  \draw[->,grijs] (-2.4,0) -- (2.5,0); \draw[->,grijs] (0,-4.3) -- (0,4.6);
  \kromme{-2.02}{2.02}{-(\x)^4+3*(\x)^2}
\end{tikzpicture}\\ \small $-x^4 + 3x^2$: links en rechts onder\end{tabular} & \begin{tabular}{c}\begin{tikzpicture}[x=.55cm,y=.25cm]
  \draw[->,grijs] (-2.4,0) -- (2.5,0); \draw[->,grijs] (0,-4.3) -- (0,4.6);
  \kromme{-2.22}{2.22}{-(\x)^3+3*\x}
\end{tikzpicture}\\ \small $-x^3 + 3x$: links boven, rechts onder\end{tabular} \\ \hline
\end{tabular}
\end{center}

Bij een even graad gaan de twee uiteinden dezelfde kant op, bij een oneven graad in tegengestelde richting. Daaruit volgt ook dat een veeltermfunctie van oneven graad altijd een nulwaarde heeft: links en rechts hebben de functiewaarden een ander teken, en de grafiek is een doorlopende lijn, dus ze moet ergens de $x$-as kruisen.

\voorbeeld{}{Hoe gedraagt de grafiek van $f(x) = (x + 1)(2 - x)^3$ zich aan de uiteinden?\label{vb:h1-uiteinden}}{%
\naast{%
\uitleg{Bepaal de hoogstegraadsterm zonder uit te werken: vermenigvuldig de hoogstegraadstermen van de factoren.}
\begin{stappenvrij}{4.3cm}{3.6cm}
x \cdot (-x)^3 = -x^4 & In $2 - x$ is de hoogste term $-x$. \\[1ex]
n = 4 \text{ even},\ a_4 = -1 < 0 & \\[1ex]
\text{links en rechts onder} & Zoals $-x^4$.
\end{stappenvrij}
Wie $(2 - x)^3$ leest als $x^3$, vindt $x^4$ en tekent de uiteinden verkeerd.}{%
\begin{tikzpicture}[x=.6cm,y=.2cm,font=\scriptsize]
  \rooster{-2}{1}{3}{-6}{2}{8}
  \kromme{-1.15}{3}{(\x+1)*(2-\x)^3}
  \vol{(-1,0)} \vol{(2,0)}
\end{tikzpicture}}}

\begin{oefeningen}
\needspace{6\baselineskip}
\oef{h1-uiteinden}
\end{oefeningen}

\subsection{Snijden of raken}

In een nulwaarde kan de grafiek door de $x$-as gaan of ze alleen raken. De figuren tonen $(x - 1)^m$ voor $m = 1, 2, 3$: de multipliciteit beslist.

\figuren{%
\begin{tikzpicture}[x=.8cm,y=.8cm,font=\scriptsize]
  \draw[->,grijs] (-.4,0) -- (2.5,0) node[right,text=black] {$x$}; \draw[->,grijs] (0,-1.3) -- (0,1.5);
  \kromme{-.2}{2.2}{\x-1} \vol{(1,0)} \node[below] at (1,-1.4) {$x - 1$: snijdt};
\end{tikzpicture}\qquad
\begin{tikzpicture}[x=.8cm,y=.8cm,font=\scriptsize]
  \draw[->,grijs] (-.4,0) -- (2.5,0) node[right,text=black] {$x$}; \draw[->,grijs] (0,-1.3) -- (0,1.5);
  \kromme{-.15}{2.15}{(\x-1)^2} \vol{(1,0)} \node[below] at (1,-1.4) {$(x - 1)^2$: raakt};
\end{tikzpicture}\qquad
\begin{tikzpicture}[x=.8cm,y=.8cm,font=\scriptsize]
  \draw[->,grijs] (-.4,0) -- (2.5,0) node[right,text=black] {$x$}; \draw[->,grijs] (0,-1.3) -- (0,1.5);
  \kromme{-.05}{2.08}{(\x-1)^3} \vol{(1,0)} \node[below] at (1,-1.4) {$(x - 1)^3$: snijdt};
\end{tikzpicture}}

\begin{theorie}
\begin{stelling}{Teken rond een nulwaarde}
Is $a$ een nulwaarde met multipliciteit $m$ van een veeltermfunctie $f$, dan

\qquad bij $m$ \textbf{\textcolor{accent}{oneven}}: \ wisselt het teken van $f$ in $a$, en \textbf{\textcolor{accent}{snijdt}} de grafiek de $x$-as;

\qquad bij $m$ \textbf{\textcolor{accent}{even}}: \ blijft het teken van $f$ hetzelfde, en \textbf{\textcolor{accent}{raakt}} de grafiek de $x$-as.
\end{stelling}

\begin{bewijs}
Schrijf $f(x) = (x - a)^m \cdot Q(x)$ met $Q(a) \neq 0$. De grafiek van $Q$ maakt geen sprongen, dus dicht bij $a$ heeft $Q(x)$ hetzelfde teken als $Q(a)$, links en rechts van $a$. Het teken van $f$ verandert dus alleen door de factor $(x - a)^m$. Rechts van $a$ is $x - a > 0$, dus $(x - a)^m > 0$. Links van $a$ is $x - a < 0$, en dan is
\bewijsregel{(x - a)^m > 0 \text{ als } m \text{ even is}, \qquad (x - a)^m < 0 \text{ als } m \text{ oneven is.}}
Bij een oneven $m$ heeft $f$ dus links en rechts van $a$ een ander teken, bij een even $m$ hetzelfde.
\end{bewijs}
\end{theorie}

Zo stel je een tekentabel snel op. Rechts van de grootste nulwaarde heeft $f$ het teken van $a_n$ (gedrag aan de uiteinden). Ga dan naar links: bij elke nulwaarde met oneven multipliciteit wisselt het teken, bij een even multipliciteit niet.

\voorbeeld{}{Stel de tekentabel op van $f(x) = (x + 2)^3(x - 1)^2(x - 3)$.\label{vb:h1-teken}}{%
\uitleg{Je hoeft geen getallen in te vullen. Begin rechts met het teken van de hoogstegraadsterm, en ga van rechts naar links.}
\begin{stappen}
x^3 \cdot x^2 \cdot x = x^6 & Rechts van $3$ is $f$ positief. \\[1ex]
3:\ \text{enkel} & Het teken wisselt: tussen $1$ en $3$ negatief. \\[1ex]
1:\ \text{dubbel} & Het teken blijft: tussen $-2$ en $1$ negatief. \\[1ex]
-2:\ \text{multipliciteit } 3 & Het teken wisselt: links van $-2$ positief.
\end{stappen}
\figuren{\tekentabel{cccccccc}{x & & -2 & & 1 & & 3 & }{f(x) & + & 0 & - & 0 & - & 0 & +}}}

\begin{werkwijze}
Ontbind volledig en noteer de nulwaarden met hun multipliciteit. Zet ze geordend in de tekentabel. Begin rechts met het teken van $a_n$, en ga naar links: wissel van teken bij een oneven multipliciteit, behoud het bij een even multipliciteit. Factoren zonder nulwaarden, zoals $x^2 + 1$, veranderen het teken nergens.
\end{werkwijze}

\begin{oefeningen}
\needspace{6\baselineskip}
\oef{h1-teken}
\end{oefeningen}

\subsection{De grafiek schetsen}

Nu heb je alles om de grafiek van een veeltermfunctie te schetsen: de nulwaarden en hun multipliciteit, de tekentabel, het gedrag aan de uiteinden en het snijpunt met de $y$-as.

\voorbeeld{}{Schets de grafiek van $f(x) = -(x + 1)^2(x - 2)$.\label{vb:h1-schets}}{%
\naast{%
\uitleg{Lees de nulwaarden en hun multipliciteit af, en het gedrag aan de uiteinden. Bereken het snijpunt met de $y$-as.}
\begin{stappenvrij}{4.5cm}{3.3cm}
\text{hoogstegraadsterm } -x^3 & Links boven, rechts onder. \\
-1 \text{ (dubbel)},\ 2 \text{ (enkel)} & Raakt in $-1$, snijdt in $2$. \\
f(0) = -1 \cdot (-2) = 2 & Snijpunt met de $y$-as.
\end{stappenvrij}
\figuren{\tekentabel{cccccc}{x & & -1 & & 2 & }{f(x) & + & 0 & + & 0 & -}}}{%
\begin{tikzpicture}[x=.6cm,y=.42cm,font=\scriptsize]
  \rooster{-3}{1}{3}{-4}{1}{5}
  \kromme{-2.1}{2.3}{-(\x+1)^2*(\x-2)}
  \vol{(-1,0)} \vol{(2,0)} \vol{(0,2)}
\end{tikzpicture}}}

\begin{werkwijze}
Ontbind en bepaal de nulwaarden met hun multipliciteit. Stel de tekentabel op. Bepaal het gedrag aan de uiteinden met de hoogstegraadsterm en bereken $f(0)$. Teken: snijden bij oneven, raken bij even multipliciteit. Waar precies de toppen liggen, bereken je later met afgeleiden.
\end{werkwijze}

\begin{oefeningen}
\oef{h1-schets}
\end{oefeningen}

Is het voorschrift uitgewerkt gegeven, dan moet je eerst ontbinden. Dat is meestal het meeste werk.

\voorbeeld{}{Schets de grafiek van $f(x) = x^4 - 2x^3 - 3x^2 + 4x + 4$.\label{vb:h1-schets4}}{%
\naast{%
\uitleg{Zoek nulwaarden bij de delers van $4$. Deel met Horner, en probeer dezelfde nulwaarde opnieuw op het quotiënt.}
\begin{stappenvrij}{5.1cm}{3.6cm}
f(-1) = 1 + 2 - 3 - 4 + 4 = 0 & $x + 1$ is een deler. \\
f(x) = (x + 1)(x^3 - 3x^2 + 4) & Horner: $1, -3, 0, 4$. \\
= (x + 1)^2(x^2 - 4x + 4) & Nog eens Horner met $-1$. \\
= (x + 1)^2(x - 2)^2 & Merkwaardig product. \\
\text{hoogstegraadsterm } x^4 & Links en rechts boven. \\
f(0) = 4 & Snijpunt met de $y$-as.
\end{stappenvrij}
Twee dubbele nulwaarden: de grafiek raakt de $x$-as twee keer en komt nooit onder de as, want $f(x)$ is een kwadraat.}{%
\begin{tikzpicture}[x=.58cm,y=.25cm,font=\scriptsize]
  \rooster{-2}{1}{3}{0}{2}{10}
  \kromme{-1.8}{2.8}{(\x+1)^2*(\x-2)^2}
  \vol{(-1,0)} \vol{(2,0)} \vol{(0,4)}
\end{tikzpicture}}}

\begin{oefeningen}
\needspace{6\baselineskip}
\oef{h1-schets2}
\end{oefeningen}

%========================================================================================
\opdrachtenkop

\begin{oefeningen}
\needspace{12\baselineskip}
\oef{h1-o1}

\oef{h1-o13}

\needspace{12\baselineskip}
\oef{h1-o3}

\oef{h1-o10}

\needspace{12\baselineskip}
\oef{h1-o15}

\oef{h1-o2}

\needspace{8\baselineskip}
\oef{h1-o12}

\needspace{8\baselineskip}
\oef{h1-o14}

\oef{h1-o11}

\oef{h1-o4}

\needspace{8\baselineskip}
\oef{h1-o17}

\needspace{14\baselineskip}
\oef{h1-o19}

\needspace{9\baselineskip}
\oef{h1-o20}

\oef{h1-o6}

\needspace{8\baselineskip}
\oef{h1-o16}

\needspace{10\baselineskip}
\oef{h1-o23}

\oef{h1-o7}

\needspace{8\baselineskip}
\oef{h1-o21}

\oef{h1-o5}

\needspace{8\baselineskip}
\oef{h1-o22}

\needspace{14\baselineskip}
\oef{h1-o18}

\oef{h1-o8}

\needspace{8\baselineskip}
\oef{h1-o24}

\needspace{10\baselineskip}
\oef{h1-o25}

\needspace{8\baselineskip}
\oef{h1-o27}

\needspace{8\baselineskip}
\oef{h1-o26}

\needspace{6\baselineskip}
\oef{h1-o9}
\end{oefeningen}

%----------------------------------------------------------------------------------------
% KOPPEL: verbanden met andere leerstof
%----------------------------------------------------------------------------------------
\needspace{20\baselineskip}
\opdrachtsectie{Koppel}

\lkr{Algemene functies: kijk na wat al gezien is. Het pakket legt de volgorde niet vast.}

\begin{oefeningen}

\oef{h1-koppel-3}
\end{oefeningen}

%----------------------------------------------------------------------------------------
\opdrachtsectie{Test jezelf}

\begin{oefeningen}
\oef{h1-t1}

\oef{h1-t2}

\oef{h1-t5}

\oef{h1-t3}

\oef{h1-t4}

\oef{h1-t6}
\end{oefeningen}

\herhaalschema{%
  \fv{oef:h1-t1}{vb:h1-horner}{oef:h1-horner}%
  \fv{oef:h1-t2}{vb:h1-staartdeling}{oef:h1-staart}%
  \fv{oef:h1-t5}{vb:h1-factor}{oef:h1-factor}%
  \fv{oef:h1-t3}{vb:h1-ontbinden}{oef:h1-ontbinden}%
  \fv{oef:h1-t4}{vb:h1-schets}{oef:h1-schets}%
  \fv{oef:h1-t6}{vb:h1-voorschrift}{oef:h1-voorschrift}}

\ifzelfstudie
\section*{Maak je eigen samenvatting}
Sluit dit boek en schrijf in je schrift: de graad van een product, de euclidische deling, het schema van Horner, de reststelling en de factorstelling met hun bewijs, hoe je kandidaat-nulwaarden vindt en waarom, het aantal nulwaarden, het gedrag aan de uiteinden, wat multipliciteit betekent voor de grafiek, en de werkwijze om een veeltermfunctie te schetsen. Vergelijk daarna met de kaders.
\fi
